\documentclass[10pt]{moderncv}

\moderncvstyle{classic}

\usepackage{amsmath, amssymb}
\usepackage{hyperref}

\firstname{David}
\familyname{Roe}
\title{}
\address{Department of Mathematics \\ Massachusetts Institute of Technology \\ 2-336, 77 Massachusetts Ave}{Cambridge, MA 02139}
\email{roed@mit.edu}
\extrainfo{\weblink{https://math.mit.edu/\textasciitilde roed}}

\newcommand{\goodindentsection}[3]{\section{#1} \cvline{}{} \cvline[-12pt]{#2}{#3}}
\newcommand{\goodindententrysection}[7]{\section{#1} \cvline{}{} \cvline[-12pt]{#2}{
    {\bfseries#3}%
    \ifthenelse{\equal{#4}{}}{}{, {\slshape#4}}%
    \ifthenelse{\equal{#5}{}}{}{, #5}%
    \ifthenelse{\equal{#6}{}}{}{, #6}%
    .%
    \ifthenelse{\equal{#7}{}}{}{\newline{}\small#7}}}
\newcommand{\cvlinesm}[2]{\cvline[-8pt]{#1}{#2}}
\newcommand{\bluelink}[2]{\httplink[\color{blue}{#1}]{#2}}

\newcommand{\ZZ}{\mathbb{Z}}
\newcommand{\CC}{\mathbb{C}}
\DeclareMathOperator{\PGL}{PGL}

\begin{document}


\maketitle
%% Education %%
\section{\textsc{Education}}
\cventry{2011}{Ph.D. in Mathematics}{Harvard University}{Cambridge, MA}{}{}
\cvline{}{\textsc{Ph.D. Thesis}}
\cvline{title}{The Local Langlands Correspondence for Tamely Ramified Groups}
\cvline{supervisor}{Professor Benedict Gross}

\cventry{2006}{B.S. in Mathematics and in Literature}{Massachusetts Institute of Technology}{Cambridge, MA}{}{}

%% Employment %%
\goodindentsection{\textsc{Employment}}{2018--present}{Research Scientist at MIT}

\cvlinesm{2015--2017}{Postdoctoral Fellow at the University of Pittsburgh (with Thomas Hales)}

\cvlinesm{2014--2015}{Postdoctoral Fellow at the University of British Columbia (with Julia Gordon)}

\cvlinesm{2011--2014}{Postdoctoral Fellow at the University of Calgary (with Clifton Cunningham)}

%% Research Interests %%
%\section{\textsc{Research Interests}}

%\cvline{}{Computational number theory, arithmetic geometry, mathematical databases, $p$-adics}

%% Editorial Positions %%
\goodindententrysection{\textsc{Editorial Positions}}{Spring 2024}{Program committee}{ANTS: Algorithmic number theory symposium}{}{}{\url{antsmath.org}}

\cventry{Spring 2023}{Associate editor}{LuCaNT: LMFDB, Computation, and Number Theory}{}{}{\url{lucant.org}}

\cventry{2020--2023}{Associate editor}{LMFDB: L-functions and modular forms database}{}{}{\url{lmfdb.org}}

%% Honors and Awards %%
\goodindentsection{\textsc{Honors and Awards}}{2022}{MIT award for service to the Math Community}

\cvlinesm{2020}{2020 COVID-19 Hero Award for creating researchseminars.org}

\cvlinesm{2011}{PIMS Postdoctoral Fellow}

\cvlinesm{2006}{National Science Foundation: Graduate Research Fellow}

\cvlinesm{2006}{National Defense Science and Engineering Graduate Research Fellow}

\cvlinesm{2006}{Phi Beta Kappa member}

\cvlinesm{2006}{Todd Anderson Award for excellence in teaching at the MIT Experimental Study Group}

\cvlinesm{2004}{Summer Program in Undergraduate Research at MIT: Rogers Prize}

\cvlinesm{2002}{National Merit Scholar}

% Grants
\goodindentsection{\textsc{Grants}}{2025--2030}{Scientific Software Research Faculty Award, Simons Foundation \newline \$250,000 + 50\% salary}

% UROPs
\goodindentsection{\textsc{UROPs supervised}}{Summer 2023}{Misheel Otgonbayar, \emph{Mod $m$ Galois representations of elliptic curves}}

\cvlinesm{Summer 2022}{Andrei Marginean, \emph{Finite subgroups of $\PGL(n, \CC)$}}

\cvlinesm{Summer 2022}{Brian Reinhart, \emph{Minimal Faithful Representations of Finite Groups}}

% Service (mathcamp, researchseminars, SageMath Code of Conduct Committee)

%% Publications %%
\section{\textsc{Preprints}}

\cvline{}{Lewis Combes, John Jones, Jennifer Paulhus, David Roe, Manami Roy, and Sam Schiavone, \bluelink{Creating a Dynamic Database of Finite Groups}{math.mit.edu/~roed/writings/research/FiniteGroups.pdf}.}

\cvline{}{Edgar Costa, Kiran Kedlaya and David Roe, \bluelink{Hypergeometric L-functions in average polynomial time II}{https://arxiv.org/abs/2310.06971}, arXiv:2310.06971, accepted to ANTS XVI.}

\section{\textsc{Publications}}

\cvline{}{Taylor Dupuy, Kiran Kedlaya, David Roe and Christelle Vincent, \bluelink{Counterexamples to a Conjecture of Ahmadi and Shparlinski}{doi.org/10.1080/10586458.2021.1980463}, Experimental Mathematics \textbf{23} (2023), no. 3, pp. 540--544.}

\cvline{}{Clifton Cunningham and David Roe, \bluelink{Commutative character sheaves and geometric types for supercuspidal representations}{doi.org/10.5802/ahl.105}, Annales Henri Lebesgue \textbf{4} (2021), pp. 1389--1420.}

\cvline{}{Alex Best, Jonathan Bober, Andrew Booker, Edgar Costa, John Cremona, Maarten Derickx, David Lowry-Duda, Min Lee, David Roe, Andrew Sutherland, John Voight, \bluelink{Computing classical modular forms}{doi.org/10.1007/978-3-030-80914-0_4}, Arithmetic Geometry, Number Theory, and Computation, Simons Symposia, Springer, 2021.}

\cvline{}{Taylor Dupuy, Kiran Kedlaya, David Roe and Christelle Vincent, \bluelink{Isogeny classes of abelian varieties over finite fields in the LMFDB}{doi.org/10.1007/978-3-030-80914-0_13}, Arithmetic Geometry, Number Theory, and Computation, Simons Symposia, Springer, 2021.}

\cvline{}{Edgar Costa and David Roe, \bluelink{Zen and the art of database maintenance}{doi.org/10.1007/978-3-030-80914-0_8}, Arithmetic Geometry, Number Theory, and Computation, Simons Symposia, Springer, 2021.}

\cvline{}{Edgar Costa, Kiran Kedlaya and David Roe, \bluelink{Hypergeometric L-functions in average polynomial time}{doi.org/10.2140/obs.2020.4.143}, The Open Book Series \textbf{4}, Fourteenth Algorithmic Number Theory Symposium, 2020.}

\cvline{}{David Roe, \bluelink{The inverse Galois problem for $p$-adic fields}{doi.org/10.2140/obs.2019.2.393}, Proceedings of the 13th Algorithmic Number Theory Symposium (ANTS-XIII).   Open Book Series \textbf{2}, Math. Sci. Pub., Berkeley CA, 2019, pp 393--409.}

\cvline{}{Clifton Cunningham and David Roe, \bluelink{From the function-sheaf dictionary to quasicharacters of $p$-adic tori}{doi.org/10.1017/S1474748015000286}, J. Inst. Math. Jussieu \textbf{17} (2018), no. 1, pp. 1--37.}

\cvline{}{Xavier Caruso, David Roe and Tristan Vaccon, \bluelink{\texttt{ZpL}: a $p$-adic precision package}{doi.org/10.1145/3208976.3208995}, Proceedings of the 2018 ACM on International Symposium on Symbolic and Algebraic Computation.  ACM, New York, 2018, pp 119--126.}

\cvline{}{Xavier Caruso, David Roe and Tristan Vaccon, \bluelink{Characteristic polynomials of $p$-adic matrices}{doi.org/10.1145/3087604.3087618}, Proceedings of the 2017 ACM on International Symposium on Symbolic and Algebraic Computation.  ACM, New York, 2017, pp 389--396.}

\cvline{}{Xavier Caruso, David Roe and Tristan Vaccon, \bluelink{Division and slope factorization of $p$-adic polynomials}{doi.org/10.1145/2930889.2930898}, Proceedings of the 2016 ICM on International Symposium on Symbolic and Algebraic Computation. ACM, New York, 2016, pp. 159--166.}

\cvline{}{Moshe Adrian and David Roe, \bluelink{Rectifiers and the local Langlands correspondence: the unramified case}{doi.org/10.4310/MRL.2016.v23.n3.a1}, Math. Res. Lett. \textbf{23} (2016), no. 3, pp. 593--619.}

\cvline{}{Xavier Caruso, David Roe and Tristan Vaccon, \bluelink{$p$-adic stability in linear algebra}{doi.org/10.1145/2755996.2756655}, Proceedings of the 2015 ICM on International Symposium on Symbolic and Algebraic Computation. ACM, New York, 2015, pp. 101--108.}

\cvline{}{Julia Gordon and David Roe, \bluelink{The canonical measure on a reductive $p$-adic group is motivic}{doi.org/10.24033/asens.2322}, Ann. Sci. \'Ec. Norm. Sup\'er. \textbf{50} (2015), no. 2, pp. 345--355.}

\cvline{}{David Roe, \bluelink{The 3-adic eigencurve at the boundary of weight space}{doi.org/10.1142/S1793042114500560}, Int. J. Number Theory \textbf{10} (2014), no. 7, pp. 1791--1806.}

\cvline{}{Xavier Caruso, David Roe and Tristan Vaccon, \bluelink{Tracking $p$-adic precision}{doi.org/10.1112/S1461157014000357}, LMS J. Comput. Math. \textbf{17} (Special issue A) (2014), 274--294.}

\cvline{}{David Roe, \bluelink{Constructing local $L$-packets for tame unitary groups$^*$}{arxiv.org/abs/1311.7456}, arXiv:1311.7456.}

\cvline{}{David Roe, \bluelink{The local Langlands correspondence for tamely ramified groups}{math.mit.edu/~roed/writings/thesis.pdf}, Ph.D. thesis, Harvard University, 2011.}

\cvline{}{Timothy G. Abbott, Kiran Kedlaya and David Roe, \bluelink{Bounding Picard numbers of surfaces using $p$-adic cohomology}{smf.emath.fr/publications/estimation-du-nombre-de-picard-des-surfaces-par-la-cohomologie-p-adique}, in \textit{Arithmetic, Geometry and Coding Theory (AGCT 2005), S\'eminaires et Congr\`es 21, Societ\'e Math\'ematique de France}, 2009, 125--159.}

%% Software %%

% goodindententrysection was too fragile for hrefs, so we break it out
\section{\textsc{Software and Databases}}
\cvline{}{}
\cvline[-12pt]{2008--2023}{\textbf{L-functions and Modular Forms Database (LMFDB)}, \bluelink{www.lmfdb.org}{lmfdb.org}.{\newline{}\small Serve on editorial board, develop site-wide infrastructure, add content (described below)}}

\cventry{2016--2023}{LMFDB: Abelian varieties over finite fields}{}{\bluelink{www.lmfdb.org/Variety/Abelian/Fq}{lmfdb.org/Variety/Abelian/Fq}}{}{}

\cventry{2018--2021}{LMFDB: Classical modular forms}{}{\bluelink{www.lmfdb.org/ModularForm/GL2/Q/holomorphic}{lmfdb.org/ModularForm/GL2/Q/holomorphic/}}{}{}

\cventry{2019--2023}{LMFDB: Finite groups}{}{\bluelink{beta.lmfdb.org/Groups/Abstract}{beta.lmfdb.org/Groups/Abstract}}{}{}

\cventry{2020--2023}{LMFDB: Hypergeometric motives}{}{\bluelink{beta.lmfdb.org/Motive/Hypergeometric/Q/}{beta.lmfdb.org/Motive/Hypergeometric/Q/}}{}{}

\cventry{2022--2023}{LMFDB: Modular curves}{}{\bluelink{beta.lmfdb.org/ModularCurve/Q/}{beta.lmfdb.org/ModularCurve/Q/}}{}{}

\cventry{2006--2023}{SageMath}{}{\bluelink{sagemath.org}{sagemath.org}}{}{Most of my effort has gone to creating $p$-adics module, but also worked on finite fields, elliptic curves, number fields, the transition to git and github, and the coercion system.}

\cventry{2020--2022}{Research Seminars}{}{\bluelink{researchseminars.org}{researchseminars.org}}{}{Cofounder of online index of online seminars worldwide}

\cventry{2021--2022}{Pset Partners}{}{\bluelink{psetpartners.mit.edu}{psetpartners.mit.edu/about}}{}{Helped with initial setup and matching algorithm}

\cventry{2019--2022}{Psycodict}{}{\bluelink{github.com/roed314/psycodict}{github.com/roed314/psycodict}}{}{Python-SQL interface, used in LMFDB, researchseminars and pset partners}

%% Talks %%
\goodindentsection{\textsc{Invited Talks}}{Jun 2024}{Finite groups and K3 surfaces in the LMFDB (Computational geometry)}

\cvlinesm{Sep 2023}{Modular curves in the LMFDB (Modular curves and Galois representations)}

\cvlinesm{Mar 2023}{Modular curves and finite groups (COUNT)}

\cvlinesm{Jan 2023}{Modular curves and finite groups (2023 Simons meeting)}

\cvlinesm{May 2022}{The L-functions and modular forms database (Big Data in Pure Mathematics)}

\cvlinesm{Feb 2021}{Algebraic structures in Sage (Sage-Oscar Days)}

\cvlinesm{Nov 2019}{A database of $p$-adic tori (UBC number theory seminar)}

\cvlinesm{Sep 2019}{A database of $p$-adic tori ($p$-adic Langlands conference)}

\cvlinesm{Jun 2019}{A database of $p$-adic tori (CMS Summer meetings)}

\cvlinesm{Jul 2018}{Inverse Galois problem for $p$-adic fields (ANTS 13)}

\cvlinesm{Feb 2018}{Using lattice to track $p$-adic precision (Numerical methods for algebraic curves)}

\cvlinesm{Sep 2017}{How to win the lottery (Reed College colloquium)}

\cvlinesm{Aug 2017}{Introduction to $p$-adics in Sage (Sage Days 88)}

\cvlinesm{July 2017}{Introduction to $p$-adics in Sage (Sage Days 87)}

\cvlinesm{June 2017}{Makdisi's algorithm for Jacobians of $p$-adic curves (Sage Days 86.5)}

\cvlinesm{May 2017}{Computing with modular forms (UNCG Summer School)}

\cvlinesm{Feb 2017}{Modular forms and modular symbols in Sage (CLap-CLap)}

\cvlinesm{Sep 2016}{Algebraic tori and a computational inverse Galois problem (PANTS 26)}

\cvlinesm{Mar 2016}{Overconvergent modular symbols (Sage Days 71)}

\cvlinesm{Sep 2015}{Overconvergent modular symbols in Sage (RIMS conference on computer algebra)}

\cvlinesm{July 2015}{Positive slope pieces of the eigencurve via interpolation (Comp. Rep. Theory in N.T.)}

\cvlinesm{May 2015}{Positive slope pieces of the eigencurve via interpolation (Explicit Methods for Ab. Var.)}

\cvlinesm{June 2014}{Geometrizing the Langlands correspondence in mixed characteristic (CNTA XIII)}

\cvlinesm{May 2014}{A function-sheaf dictionary for tori over local fields (FRG on periods and $L$-functions)}

\cvlinesm{Feb 2014}{Numerical Methods in p-adic Linear Algebra (Lethbridge Number Theory Seminar)}

\cvlinesm{Jan 2014}{A function-sheaf dictionary for tori over local fields (AMS Session: the Langlands Program)}

\cvlinesm{Dec 2013}{Quasicharacter sheaves for tori (Quebec-Vermont Number Theory Seminar)}

\cvlinesm{Dec 2013}{An Introduction to the local Langlands correspondence (UQAM combinatorics seminar)}

\cvlinesm{June 2013}{Geometrizing characters of tori (PRIMA 2013)}

\cvlinesm{May 2013}{Geometrizing characters of tori (Alberta Number Theory Days)}

\cvlinesm{Jan 2013}{Geometrizing quasicharacters of tori (UCSC Number Theory Seminar)}

\cvlinesm{Jan 2013}{Geometrizing quasicharacters of tori (AMS Session: Witt Vectors, Lifting and Descent)}

\cvlinesm{May 2012}{The local Langlands correspondence for tamely ramified groups (Pacific Northwest Number Theory Conference)}

\cvlinesm{Apr 2012}{The local Langlands correspondence for tamely ramified groups (University of Utah)}

\cvlinesm{Mar 2012}{Precision models for arithmetic in local fields (UBC Number Theory Seminar)}

\cvlinesm{Feb 2012}{The state of $p$-adics in Sage (Sage Days 36)}

\cvlinesm{Jun 2011}{Zeta functions with $p$-adic cohomology (Geocrypt 2011)}

\cvlinesm{Jun 2011}{Precision models for arithmetic in local fields (Th\'eorie de Hodge p-adique, \'equations diff\'erentielles p-adiques et leurs applications)}

\cvlinesm{Jan 2011}{The local Langlands correspondence for tamely ramified groups (University of Calgary)}

\cvlinesm{Nov 2010}{$p$-Adics in Sage: present and future (University Rennes 1)}

\cvlinesm{Oct 2010}{The local Langlands correspondence for tamely ramified groups (CCR West)}

\cvlinesm{Oct 2010}{The local Langlands correspondence for tamely ramified groups (Quebec-Maine Number Theory Conference)}

\cvlinesm{Apr 2010}{Zeta functions with $p$-adic cohomology (Counting Points: Theory, Algorithms and Practice)}

\cvlinesm{Mar 2008}{A bound for the number of automorphisms of an arithmetic Riemann surface (group project presentation, Arizona Winter School)}

\cvlinesm{Sep 2007}{$p$-Adic arithmetic in Sage (Sage Days 5)}

\cvlinesm{Jun 2007}{$p$-Adics in Sage (Sage Days 4)}

%% Teaching Experience %%
\goodindententrysection{\textsc{Teaching Experience}}{Fall 2020}{Instructor: Arithmetic geometry (18.782)}{Dept. of Mathematics, MIT}{}{}{Instructor, online due to covid}

\cventry{Spring 2019}{Instructor: $p$-adic analysis}{Dept. of Mathematics, MIT}{}{}{Independent study with 3 students}

\cventry{Fall 2017}{Instructor: Intro. Abstract Algebraic Structures (430)}{Dept. of Mathematics, Pitt}{}{}{Sole instructor.}

\cventry{Fall 2017}{Lecturer: Calculus 1 (220)}{Dept of Mathematics, Pitt}{}{}{}

\cventry{Spring 2017}{Lecturer: Calculus 3 (240)}{Dept. of Mathematics, Pitt}{}{}{}

\cventry{Fall 2016}{Lecturer: Calculus 1 (220)}{Dept. of Mathematics, Pitt}{}{}{}

\cventry{Fall 2016}{Lecturer: Discrete Mathematical Structures (400)}{Dept. of Mathematics, Pitt}{}{}{}

\cventry{Spring 2016}{Instructor: Intro. Abstract Algebraic Structures (430)}{Dept. of Mathematics, Pitt}{}{}{Sole instructor.}

\cventry{Fall 2015}{Lecturer: Calculus 1 (220)}{Dept. of Mathematics, Pitt}{}{}{Lectured for two sections.}

\cventry{Fall 2014}{Lecturer: Multivariable Calculus (200)}{Dept. of Mathematics, UBC}{}{}{Lectured for two sections.}

\cventry{Fall 2013}{Lecturer: Linear Methods (211)}{Dept. of Mathematics, Calgary}{}{}{Lectured for a 140-person section of science majors.}

\cventry{Fall 2012}{Instructor: Computational Number Theory (627)}{Dept. of Mathematics, Calgary}{}{}{Designed and taught this joint undergraduate/graduate course.}

\cventry{Fall 2008}{Course Assistant: Algebraic Number Theory (223a)}{Dept. of Mathematics, Harvard}{}{}{Taught recitations, assisted students in preparing for weekly student talks, and graded.}

\cventry{Fall 2007, 2008}{Teaching Fellow: Magic of Numbers (QR28)}{Core Program, Harvard}{}{}{Taught recitations, wrote homework and tests, and graded.}

\cventry{Summers 2007-2009 \\ 2011 \\ 2016}{Mentor and Academic Coordinator}{Canada/USA Mathcamp}{}{}{Taught short courses at this camp for high school students, including Unique Factorization, Algebraic Topology, $p$-Adics, Elliptic Curves, Modular Forms, Coxeter Groups, Factoring Algorithms.} 

\cventry{Spring 2006}{Teaching Assistant: Project Lab in Math. (18.821)}{Dept. of Mathematics, MIT}{}{}{Assisted small groups of students with carrying out and writing up a series of research projects.}

\cventry{Fall 2005 \\ Spring 2005 \\ Fall 2003}{Instructor: Multivariable Calculus (18.02)}{Experimental Study Group, MIT}{}{}{Sole instructor.  Lectured, wrote homework and tests, and graded.}

%% Conferences Organized %%
\goodindentsection{\textsc{Conferences Organized}}{April 2024}{\bluelink{Hypergeometric motives workshop}{edgarcosta.org/HGMworkshop24} (MIT)}

\cvlinesm{March 2024}{\bluelink{Modular curves workshop 3}{https://math.mit.edu/~edgarc/MCW3.html} (MIT)}

\cvlinesm{November 2022}{\bluelink{Modular curves workshop 2}{https://math.mit.edu/~edgarc/MCW2.html} (MIT)}

\cvlinesm{April 2022}{\bluelink{Explicit methods in modularity}{https://researchseminars.org/seminar/ExplicitModularity} (online)}

\cvlinesm{March 2022}{\bluelink{Modular curves workshop}{https://math.mit.edu/~edgarc/MCW.html} (MIT)}

\cvlinesm{July 2017}{\bluelink{Sage Days 87: $p$-adics in Sage}{https://wiki.sagemath.org/days87} (University of Vermont)}

\cvlinesm{Mar 2016}{\bluelink{Explicit $p$-adic Methods in Number Theory}{http://wiki.sagemath.org/days71} (Oxford)}

\cvlinesm{Dec 2015}{\bluelink{CMS Winter Meeting: Representation Theory session}{https://math.mit.edu/~roed/conferences/CMS15/} (Montreal, QC)}

\cvlinesm{May 2013}{\bluelink{Alberta Number Theory Days}{https://math.mit.edu/~roed/conferences/ANTD13/} (Banff, AB)}

\cvlinesm{Feb 2013}{\bluelink{Sage Days 44: $p$-adic overconvergent symbols in Sage}{http://wiki.sagemath.org/sagedays44} (University of Wisconsin)}

\cvlinesm{March 2012}{\bluelink{Sage Days 40.5: Bug fixing workshop}{http://wiki.sagemath.org/bug19} (Seattle, WA)}

\cvlinesm{Feb 2012}{\bluelink{Sage Days 36: $p$-adics in Sage}{http://wiki.sagemath.org/padicSageDays} (University of California, San Diego)}

%\vspace{0.4in}

%% Sage %%

%\goodindentsection{\textsc{Contributions to Sage}}{}{Attended 21 Sage Days workshops, developed Sage's $p$-adic codebase (2006-2023)}
%\cvlinesm{}{Added more than 40,000 lines (net) of peer-reviewed code to the math software system Sage, including the following selected tickets viewable on \weblink{http://trac.sagemath.org}}
%\cvlinesm{\weblink[\#285]{http://trac.sagemath.org/ticket/285}}{Major new version of $p$-adic arithmetic}
%\cvlinesm{\weblink[\#721]{http://trac.sagemath.org/ticket/721}}{Tate's algorithm for elliptic curves over number fields (based on code of Cremona)}
%\cvlinesm{\weblink[\#812]{http://trac.sagemath.org/ticket/812}}{Overconvergent modular symbols (with others)}
%\cvlinesm{\weblink[\#7931]{http://trac.sagemath.org/ticket/7931}}{Improved $k$-th roots for finite fields and $\ZZ/n\ZZ$}
%\cvlinesm{\weblink[\#8335]{http://trac.sagemath.org/ticket/8335}}{Lattice of compatibly embedded finite fields via Conway polynomials (with others)}
%\cvlinesm{\weblink[\#12116]{http://trac.sagemath.org/ticket/12116}}{Improved perfect power function for integers}
%\cvlinesm{\weblink[\#12415]{http://trac.sagemath.org/ticket/12415}}{Rewrite of Sage's testing framework}
%\cvlinesm{\weblink[\#12555]{http://trac.sagemath.org/ticket/12555}}{Reimplementation of $p$-adics with better extensibility}
%\cvlinesm{\weblink[\#12766]{http://trac.sagemath.org/ticket/12766}}{Better plotting of elliptic curves}
%\cvlinesm{\weblink[\#14304]{http://trac.sagemath.org/ticket/14304}}{Updated implementation of unramified extensions}
%\cvlinesm{\weblink[\#20348]{http://trac.sagemath.org/ticket/20348}}{$p$-adic floating point rings}
%\cvlinesm{\weblink[\#23218]{http://trac.sagemath.org/ticket/23218}}{General extensions of $p$-adic rings}

%\vspace{0.4in}

%\goodindentsection{\textsc{Contributions to LMFDB}}{}{LMFDB editor, attended many LMFDB workshops; selected contributions:}
%\cvlinesm{\weblink[\#865]{https://github.com/LMFDB/lmfdb/pull/865} \weblink[\#896]{https://github.com/LMFDB/lmfdb/pull/896}}{Centralized functions for parsing search boxes}
%\cvlinesm{\weblink[\#1966]{https://github.com/LMFDB/lmfdb/pull/1966} \weblink[\#3340]{https://github.com/LMFDB/lmfdb/pull/3340}}{Isogeny classes of abelian varieties over finite fields (with Dupuy, Kedlaya and Vincent)}
%\cvlinesm{\weblink[\#2557]{https://github.com/LMFDB/lmfdb/pull/2557} \weblink[\#2572]{https://github.com/LMFDB/lmfdb/pull/2572}}{Switch from MongoDB to Postgres (with Costa)}
%\cvlinesm{\weblink[\#2717]{https://github.com/LMFDB/lmfdb/pull/2717}}{Classical modular forms (with Costa, Sutherland, Voight)}
%\cvlinesm{\weblink[\#2736]{https://github.com/LMFDB/lmfdb/pull/2736}}{Improvements to knowledge db (with Costa)}

%% Service %%
\goodindentsection{\textsc{Service}}{Spring 2024}{Co-organizer for MIT number theory seminar}

\cvlinesm{2023}{Served on NSF grant review panel}

\cvlinesm{2020--2023}{Served on LMFDB editorial board}

\cvlinesm{2020}{Created \weblink{researchseminars.org}}

\cvlinesm{2017--2024}{Served on SageMath Code of Conduct Committee}

\cvlinesm{Summer 2016}{Academic Coordinator for Canada/USA Mathcamp}

\cvlinesm{Fall 2014}{Organizer of the number theory video seminar at UBC}

\cvlinesm{Fall 2012}{Organizer of the ABC seminar at University of Calgary}

\cvlinesm{2008--2011}{Organizer of the Alcove Seminar, a graduate number theory seminar at Harvard}

\cvlinesm{2002--2009}{Volunteer teacher at Splash, a program for high school students run by MIT's Educational Studies Program}

\cvlinesm{1998--2002}{Founder and coach of one MATHCOUNTS team (\weblink{www.mathcounts.org}), coach of another}

%\subsection{\textsc{Non-academic Interests}}

%\cvlinesm{}{I enjoy hiking, bridge, board games, traveling, reading, and Mystery Hunt puzzles.  I have recently started skiing and sailing.}

%% References %%
\section{\textsc{References}}

\cventry{}{Professor Clifton Cunningham}{University of Calgary}{\emaillink{cunning@math.ucalgary.ca}}{}{Postdoctoral supervisor, collaborator}

\cventry{}{Professor Julia Gordon}{University of British Columbia}{\emaillink{gor@math.ubc.ca}}{}{Postdoctoral supervisor, collaborator}

\cventry{}{Professor Benedict Gross}{Harvard University}{\emaillink{gross@math.harvard.edu}}{}{Ph.D. thesis advisor}

\cventry{}{Professor Thomas Hales}{University of Pittsburgh}{\emaillink{hales@pitt.edu}}{}{Postdoctoral supervisor}

\cventry{}{Professor Kiran Kedlaya}{University of California, San Diego}{\emaillink{kedlaya@ucsd.edu}}{}{Undergrad research supervisor, collaborator}

\cventry{}{Andrew Sutherland}{MIT}{\emaillink{drew@math.mit.edu}}{}{Research supervisor}

\cventry{}{John Voight}{Dartmouth College}{\emaillink{john.voight@dartmouth.edu }}{}{LMFDB editor, collaborator}

\end{document}